\documentclass[11pt,a4paper]{article}
\usepackage[margin=1in]{geometry}
\usepackage{hyperref}
\usepackage{enumitem}
\usepackage{graphicx}
\usepackage{xcolor}

% -----------------------------------------------------
% bvraghav-tiet-ucs503p-article.sty
% -----------------------------------------------------

\makeatletter \newcommand{\BvrSetLabInstructor}[1]{%
  \gdef\Bvr@LabInstructor{#1}}
\newcommand{\Bvr@LabInstructor}{Prof. Asif}
\newcommand{\BvrLabInstructorValue}{\Bvr@LabInstructor}
\makeatother

\usepackage{titling}
\pretitle{\begin{center}\bfseries\fontsize{18bp}{18bp}\selectfont}
\makeatletter
\newcommand{\BvrSubtitle}[1]{%
  \posttitle{%
    \par\end{center}
    \begin{center}\large#1\end{center}
    \vskip 0.5em}%
}
\makeatother

\newcommand{\BvrHint}[1]{%
  {\normalfont\normalsize\itshape\hfill #1}}

% -----------------------------------------------------

\title{Adaptive Continuous Authentication \& Account-Takeover Detection Using Behavioral Intelligence}

\BvrSetLabInstructor{Prof. Asif}

\BvrSubtitle{%
  Project Proposal for UCS503P  \\
  Submitted to: \BvrLabInstructorValue \\ \bfseries
  Thapar Institute of Engineering and Technology%
}

\author{
  Anjali Kumari
  \thanks{Roll No: \texttt{1024030457}, %
    Email: \texttt{akumari\_be24\@thapar.edu}}
  \and
  Yatharth Kansal
  \thanks{Roll No: \texttt{1024030456}, %
    Email: \texttt{ykansal\_be24\@thapar.edu}}
  \and
  Ridhi Batra
  \thanks{Roll No: \texttt{1024030464}, %
    Email: \texttt{rbatra\_be24\@thapar.edu}}
}

\date{\today}

\begin{document}
\maketitle

\section{Higher-order goal%
  \BvrHint{\ldots secondary framing}}
This project strengthens session-level trust in authenticated systems
by continuously verifying that the person behind the keyboard is
still the legitimate user, not just at login. The engineering
objective is to translate \emph{behavioral risk signals} into a
measurable, deployable trust-scoring service.

\section{Time-to-value %
\BvrHint{\ldots primary heuristic}}
\begin{itemize}[leftmargin=0.7cm]
    \item Deliver a working end-to-end pipeline early --- capture, baseline, live score --- before perfecting any one detection model.
    \item Iterate weekly, starting with a single behavioral signal (keystroke dynamics) before fusing in more.
    \item Success in the first iteration is a live 0--100 trust score for one signal; later iterations add signals and refine accuracy.
\end{itemize}

\section{Problem Statement}
Traditional authentication is a one-time gate: credentials and MFA
are checked once at login, after which the session is marked
``trusted'' and left unmonitored. This gap lets attackers exploit:
\begin{itemize}[leftmargin=0.7cm]
    \item \textbf{Session hijacking:} a stolen token inherits full trust without presenting credentials.
    \item \textbf{Unattended workstations:} a logged-in session can be used by anyone physically present.
    \item \textbf{Binary trust:} sessions are simply ``trusted'' or not, with no way to react as risk builds mid-session.
\end{itemize}
Privilege escalation and data exfiltration during an already-authenticated
session routinely bypass login-time defenses, since nothing re-checks
identity once the door is open.

\section{Proposed Solution %
\BvrHint{\ldots Software Proposition}}
\subsection{Overview}
Build a \textbf{continuous, adaptive trust engine} that fuses
behavioral, system, and contextual signals into a single live
session-trust score (0--100), and reacts automatically as it changes:
\textbf{80--100} continue normally, \textbf{60--79} monitor closely,
\textbf{40--59} step-up MFA, \textbf{$<$40} restrict/alert.
Components: client-side signal capture (keystroke, mouse, app usage),
a baseline \& fusion engine, anomaly/sequence detection, and the
trust-scoring service driving these responses.

\subsection{Core workflow}
\begin{enumerate}[leftmargin=0.7cm]
    \item User authenticates normally; the session begins with an initial trust score.
    \item A client agent streams behavioral/system events to the backend.
    \item The fusion engine compares events against the user's baseline and updates the score in real time.
    \item The system acts per the current score band (monitor, step-up MFA, or restrict/alert).
\end{enumerate}

\section{Solution Approach %
\BvrHint{\ldots Engineering focus}}
This is an engineering/applied-ML course project, so we focus on a
working, evaluable pipeline rather than novel research in behavioral
biometrics.
\begin{itemize}[leftmargin=0.7cm]
    \item \textbf{Signal modeling:} start with tractable features (keystroke dwell/flight time, mouse velocity); use lightweight statistical or one-class-classifier baselines before heavier models. Anomaly scores feed the response layer --- never auto-terminate on one weak signal.
    \item \textbf{Architecture:} a client agent (capture), a scoring backend API (baselines, scoring, responses), and a database (baselines, event history, score timelines).
    \item \textbf{CI/CD:} automated build/test on every change, repeatable staging deploys, enabling safe frequent releases as signals are added.
\end{itemize}

\section{Evaluation Criterion %
\BvrHint{\ldots measurable and attributable}}
\subsection{Primary metric}
\textbf{Detection latency:} time between an injected anomalous
behavior (simulated takeover) and the score crossing into the
step-up/restrict band. Target: a small, bounded number of
events during pilot testing, measured from event vs.\ score-change
timestamps.

\subsection{Secondary metrics}
\begin{itemize}[leftmargin=0.7cm]
    \item \textbf{False-positive / false-negative rate} of the score bands.
    \item \textbf{Baseline convergence time} for a new user.
    \item \textbf{Reliability} of the scoring service during pilot/demo sessions.
\end{itemize}

\subsection{Pilot validation plan %
\BvrHint{\ldots high-level}}
Recruit a small group of pilot users to build genuine baselines, then
simulate takeovers (a second person continuing a logged-in session)
to test detection latency and false-negative rate across iterations.

\section{Scalability %
\BvrHint{\ldots with foresight}}
\begin{itemize}[leftmargin=0.7cm]
    \item \textbf{Theoretically:} decouple signal capture from scoring, process events via a queue/stream, and keep the scoring API stateless where feasible.
    \item \textbf{Practically:} run on common infrastructure --- managed database, containerized/platform-hosted deployment, CI/CD for staging/production.
\end{itemize}

\section{Engine availability heuristic}
Mature tooling covers this project's needs: a web/backend framework
for the API and dashboard, a managed database for event/session
storage, classical ML libraries for baselines and anomaly detection,
browser/desktop APIs for keystroke and mouse capture, a test
framework, and a CI runner with staging target. We use these to focus
on signal design, fusion logic, and evaluation.

\section{Project scope and deliverables %
\BvrHint{\ldots iteration-friendly}}
\subsection{Initial deliverable}
\begin{itemize}[leftmargin=0.7cm]
    \item Client-side capture of keystroke dynamics; baseline for a single test user
    \item Real-time trust-score computation (0--100) with a minimal live dashboard
    \item CI: build + backend unit tests + linting; CD to staging
\end{itemize}

\subsection{Subsequent deliverables}
\begin{itemize}[leftmargin=0.7cm]
    \item Mouse dynamics and app/process usage signals, fused into one score
    \item Adaptive response tiers wired to real actions (step-up MFA, session restriction)
    \item Anomaly vs.\ attack-chain (sequence-level) detection
    \item Pilot evaluation dashboard for detection latency and error rates
\end{itemize}

\section{Risks and mitigations}
\begin{itemize}[leftmargin=0.7cm]
    \item \textbf{Privacy/consent:} collect only what is needed, anonymize where possible, get informed pilot-user consent.
    \item \textbf{Baseline cold-start:} use a conservative initial score and a defined learning period for new users.
    \item \textbf{False positives frustrating users:} keep step-up MFA lightweight; tune thresholds from pilot feedback.
    \item \textbf{Measurement ambiguity:} define score bands and log all timestamps needed for detection-latency measurement before development begins.
\end{itemize}

\section{Summary}
This project proposes a continuous, adaptive authentication system
that scores session trust in real time from behavioral, system, and
contextual signals, reacting through tiers from passive monitoring to
session restriction. It emphasizes time-to-value through
signal-by-signal iteration, CI/CD for safe frequent releases, and
measurable evaluation (detection latency, false-positive/negative
rates) --- a working, evaluable pipeline rather than research-grade
novelty in behavioral biometrics.

\end{document}
